Der völkische Kampfbegriff \emph{Antisemitismus} wird heute verwendet, um jegliche Kritik an Israel oder am Zionismus zu diffamieren. Das ist eine Begriffsverdrehung.

\begin{enumerate}
    \item \textbf{Etymologische Klarstellung:} Der Begriff \emph{Semit} bezeichnet eine Sprachfamilie – Hebräisch, Arabisch, Aramäisch, Amharisch. Juden sind nur eine Teilmenge der Semiten. Der Begriff \emph{Antisemitismus} bedeutet wörtlich \emph{Feindschaft gegen alle Semiten} – also auch gegen Araber, Äthiopier, Assyrer. Er ist eine völkische Erfindung, keine wissenschaftliche Kategorie.
    
    \item \textbf{Historische Erfindung:} Der Begriff wurde 1879 von Wilhelm Marr erfunden – nicht als wissenschaftliche Kategorie, sondern als politische Kampfvokabel, um Judenhass zu legitimieren. Er wurde von den Zionisten später umgepolt und gegen Kritiker des Zionismus eingesetzt.
    
    \item \textbf{Judenhass ist real:} Die Feindschaft gegen Juden als religiöse oder ethnische Gruppe existiert – sie muss bekämpft werden. Aber sie ist nicht dasselbe wie Kritik am Staat Israel. Der Begriff \emph{Judenhass} ist der einzig korrekte Begriff für dieses Phänomen.
    
    \item \textbf{Zionismus:} Der Zionismus ist eine nationalistische, koloniale Bewegung, die einen jüdischen Staat in Palästina errichten und verteidigen will – auch auf Kosten der palästinensischen Bevölkerung. Der Zionismus ist keine Religion, sondern eine politische Ideologie.
    
    \item \textbf{Antizionismus:} Die Kritik am Zionismus und am Staat Israel ist keine Judenfeindschaft – sie ist eine politische Position, die die Rechte aller Menschen in Palästina/Israel verteidigt. Wer Antizionismus mit Judenhass gleichsetzt, verwechselt zwei völlig verschiedene Dinge.
\end{enumerate}

Die Instrumentalisierung des völkischen Kampfbegriffs \emph{Antisemitismus} dient heute dazu:
\begin{itemize}
    \item Kritik an Israels Völkerrechtsverbrechen zu unterdrücken (UN-Berichte über Verbrechen gegen Kinder im Gazastreifen)
    \item Die Linke als judenfeindlich zu brandmarken, obwohl sie lediglich antizionistisch ist
    \item Die deutsche Staatsräson zu schützen, die Israel bedingungslos unterstützt
\end{itemize}

Die Begriffsverdrehung ist perfide: Sie nimmt den realen Judenhass, den es zu bekämpfen gilt, und nutzt ihn als moralische Keule gegen politische Gegner – und vergiftet damit den gesamten Diskurs. Die Linke muss sich gegen diese Instrumentalisierung wehren: Sie verurteilt Judenhass – aber sie verurteilt auch den Zionismus. Das ist kein Widerspruch, sondern eine klare politische Position.